\section*{Lời mở đầu}

Sự phát triển nhanh của các AI Agent trong quy trình phát triển phần mềm đã đặt ra một nhu cầu mới: agent cần nắm được ngữ cảnh của dự án -- mã nguồn, tài liệu, lịch sử thay đổi và đặc tả -- để hỗ trợ công việc một cách chính xác. Trên thực tế, ở mỗi phiên làm việc, agent thường phải tự khám phá lại tri thức của dự án bằng cách đọc và tìm kiếm trực tiếp; đây là một quá trình tốn thời gian, lặp lại và không tận dụng được các mối liên hệ sẵn có giữa các thành phần của dự án. Vấn đề càng rõ khi tri thức trải trên nhiều nguồn hoặc nhiều repository, hay khi cần bắc cầu giữa mã nguồn, tài liệu và requirement.

Báo cáo này trình bày quá trình thiết kế và hiện thực Knowledge Hub, một dịch vụ đóng vai trò bộ nhớ dài hạn và nguồn cung cấp ngữ cảnh dùng chung cho các AI Agent của một team phát triển phần mềm. Đề tài xuất phát từ chính bối cảnh thực tế đó: thay vì để mỗi agent tự dò lại tri thức, hệ thống thu thập và cấu trúc hóa tri thức của một product một lần, rồi phục vụ truy xuất có cấu trúc kèm căn cứ trích dẫn cho mọi agent trong team.

Nội dung báo cáo đi từ bối cảnh và requirement, qua các phương pháp cốt lõi và kiến trúc hệ thống, tới hiện thực, đánh giá và các hướng phát triển tiếp theo; kèm theo là mã nguồn, cụm Docker và tài liệu thiết kế như những sản phẩm bàn giao của đề tài.
